\documentclass[11pt, a4paper]{article}

\usepackage[T1]{fontenc}
\usepackage[utf8]{inputenc}
\usepackage[french]{babel}
\usepackage[margin=2cm]{geometry}
\usepackage{amsmath, amssymb}
\usepackage{booktabs}
\usepackage{enumitem}
\usepackage{titlesec}
\usepackage[expansion=false]{microtype}
\usepackage{xcolor}
\usepackage{graphicx}
\usepackage{float}
\usepackage{tikz}
\usepackage[hidelinks]{hyperref}
\usepackage{tcolorbox}
\tcbuselibrary{skins, breakable}

\definecolor{accent}{HTML}{2196F3}
\definecolor{fair}{HTML}{4CAF50}
\definecolor{alertc}{HTML}{d9534f}
\definecolor{neutral}{HTML}{FF9800}

\newtcolorbox{leconbox}[1][Leçon]{
    enhanced, breakable,
    colback=fair!8, colframe=fair!70!black,
    fonttitle=\bfseries\small, title={#1},
    left=4pt, right=4pt, top=3pt, bottom=3pt,
    boxrule=0.6pt
}
\newtcolorbox{attentionbox}[1][Attention]{
    enhanced, breakable,
    colback=alertc!8, colframe=alertc!70!black,
    fonttitle=\bfseries\small, title={#1},
    left=4pt, right=4pt, top=3pt, bottom=3pt,
    boxrule=0.6pt
}
\newtcolorbox{infobox}[1][Note]{
    enhanced, breakable,
    colback=accent!8, colframe=accent!70!black,
    fonttitle=\bfseries\small, title={#1},
    left=4pt, right=4pt, top=3pt, bottom=3pt,
    boxrule=0.6pt
}

\setlength{\parskip}{0.4em}
\setlength{\parindent}{0pt}
\titlespacing*{\section}{0pt}{1em}{0.4em}
\titlespacing*{\subsection}{0pt}{0.7em}{0.3em}
\setlist{nosep, leftmargin=1.5em}

\title{\textbf{IA Responsable, German Credit Dataset\\[0.3em]
\large Rapport accompagnant \texttt{responsiveAI-german-credit\_rendu.ipynb}}}
\author{%
IA708 --- Télécom Paris, 2026 \\[0.4em]
\normalsize \textbf{Étudiants :} \\
\normalsize Julien GIMENEZ \texttt{<julien.gimenez@telecom-paris.fr>} \\
\normalsize Hugo FANCHINI \texttt{<hugo.fanchini@telecom-paris.fr>} \\
\normalsize Paul CINTRA \texttt{<paul.cintra@telecom-paris.fr>} \\
\normalsize Yimou ZHANG \texttt{<yimou.zhang@telecom-paris.fr>} \\
\normalsize Zaher HAMADEH \texttt{<zaher.hamadeh@telecom-paris.fr>}
}
\date{}

\begin{document}
\maketitle
\vspace{-1em}

\begin{abstract}
\noindent Étude d'un modèle de scoring crédit (UCI German Credit, 1000 lignes) sous l'angle de l'IA responsable.
On évalue la régression logistique baseline sur trois axes :
\textbf{équité} (Demographic Parity, Equal Opportunity, Disparate Impact, audit intersectionnel),
\textbf{interprétabilité} (SHAP global et local, détection de proxies),
et \textbf{quantification d'incertitude} (bootstrap, $B = 200$, en substitution à la robustesse adversariale sur conseil de l'enseignante).
Stack : \texttt{scikit-learn}, \texttt{fairlearn}, \texttt{shap}, exécution reproductible via \texttt{papermill}.
\end{abstract}

\tableofcontents
\vspace{1em}

% ════════════════════════════════════════════════════════════════
\section{Données et préparation}

\subsection{Dataset et attributs sensibles}

UCI Statlog German Credit (Hofmann, 1994) : 1000 lignes, 20 attributs (6 numériques, 14 catégoriels). Les codes \texttt{Axx} sont traduits en libellés français. Cible binaire : défaut de paiement (30\,\%) vs bon payeur (70\,\%).

Deux attributs sensibles sont définis :
\begin{itemize}
    \item \textbf{Genre} (depuis \texttt{personal\_status\_sex}) : 690 hommes, 310 femmes.
    \item \textbf{Âge} binarisé au seuil 25 ans, convention de la littérature (Kamiran \& Calders 2012, AIF360) : 810 \emph{older}, 190 \emph{younger}.
\end{itemize}

\subsection{Biais bruts}

Les taux de défaut diffèrent fortement entre groupes :
\begin{table}[H]
\centering
\begin{tabular}{lcc}
\toprule
\textbf{Groupe} & \textbf{P(bon)} & \textbf{P(défaut)} \\
\midrule
Homme    & 0.72 & 0.28 \\
Femme    & 0.65 & 0.35 \\
\midrule
Older    & 0.73 & 0.27 \\
Younger  & 0.58 & \textbf{0.42} \\
\bottomrule
\end{tabular}
\caption{Taux de défaut conditionnel par groupe sensible (données brutes).}
\end{table}

\noindent Un modèle qui copie ces statistiques sera mécaniquement biaisé. C'est l'enjeu central du projet.

\subsection{Préparation}

\texttt{ColumnTransformer} (Pedregosa et al.\ 2011) combine \texttt{StandardScaler} pour les numériques et \texttt{OneHotEncoder} pour les catégorielles : 56 features finales.

Split stratifié 60/20/20 (train/val/test, \texttt{random\_state=42}). Les attributs sensibles (\texttt{personal\_status\_sex}, \texttt{age\_in\_years}) sont retirés des features ; on verra que ce retrait ne suffit pas à supprimer leur influence (\S\ref{sec:proxies}).

% ════════════════════════════════════════════════════════════════
\section{Modèle baseline : régression logistique}
\label{sec:baseline}

\subsection{Choix du modèle}

\texttt{LogisticRegression} (sklearn, solver \texttt{liblinear}, $C=1$, $L_2$). Le modèle :
$$P(\text{défaut}=1 \mid x) = \sigma(w^\top x + b)$$

Trois raisons à ce choix : (1) interprétabilité native, (2) formule SHAP exacte (Lundberg \& Lee 2017, Corollary~1), (3) naturellement bien calibré sur petit dataset.

Le seuil de classification $\tau = 0.18$ est choisi pour maximiser la balanced accuracy sur la validation. Convergence : 6 itérations effectives.

\subsection{Performance globale}

\begin{table}[H]
\centering
\begin{tabular}{lc}
\toprule
\textbf{Métrique} & \textbf{Valeur (test)} \\
\midrule
AUC ROC                        & 0.785 \\
Average Precision (AP)         & 0.607 \\
Balanced Accuracy ($\tau=0.18$) & 0.669 \\
\midrule
Confusion : TN / FP / FN / TP & 66 / 74 / 8 / 52 \\
TPR (rappel défaut)            & 86.7\,\% \\
PPV (précision)                & 41.3\,\% \\
\bottomrule
\end{tabular}
\caption{Métriques de performance du baseline sur le test (200 lignes).}
\end{table}

\begin{figure}[H]
\centering
\includegraphics[width=0.95\linewidth]{../outputs/10_baseline_perf.png}
\caption{Matrice de confusion (gauche), courbe ROC (centre, AUC = 0.785), courbe Precision-Recall (droite, AP = 0.607). Au seuil 0.18, le modèle privilégie la détection des défauts (TPR = 87\,\%) au prix d'un taux de faux positifs élevé.}
\end{figure}

\subsection{Coût métier (matrice UCI)}

UCI fournit une matrice de coût : un défaut approuvé (FN) coûte 5$\times$ plus qu'un bon refusé (FP). On compare deux seuils :
\begin{itemize}
    \item \texttt{thr\_balacc} = 0.18 : maximise la balanced accuracy. Coût test = 114.
    \item \texttt{thr\_cost} = 0.14 : minimise $5 \cdot \text{FN} + 1 \cdot \text{FP}$. Coût test = 122.
\end{itemize}
Le rappel des défauts (TPR) est identique pour les deux seuils. Sur petit dataset, la calibration coût-optimale par validation se généralise mal au test.

% ════════════════════════════════════════════════════════════════
\section{Audit d'équité}

\subsection{Métriques utilisées}

Calculs via \texttt{fairlearn.metrics.MetricFrame} (Bird et al.\ 2020) :
\begin{itemize}
    \item $|\Delta\text{DP}|$ \emph{Demographic Parity gap} : $|\Pr(\hat{Y}{=}1 \mid A{=}0) - \Pr(\hat{Y}{=}1 \mid A{=}1)|$
    \item $|\Delta\text{EO}|$ \emph{Equalized Odds gap} : $\max(|\Delta\text{TPR}|, |\Delta\text{FPR}|)$ (Hardt et al.\ 2016)
    \item DI \emph{Disparate Impact ratio} : $\min_g \text{approval}_g / \max_g \text{approval}_g$. Règle US 80\,\% (EEOC) : DI $< 0.8$ = présomption de discrimination indirecte.
\end{itemize}

\begin{attentionbox}[Théorème d'impossibilité (Chouldechova 2017, Kleinberg et al.\ 2017)]
Si les taux de base diffèrent entre groupes ($p_0 \neq p_1$), DP et EO ne peuvent être satisfaits simultanément. Il faut \textbf{choisir} le critère prioritaire.
\end{attentionbox}

\subsection{Audit par axe}

\begin{table}[H]
\centering
\begin{tabular}{lcccc}
\toprule
\textbf{Attribut} & \textbf{Approval rate} (par groupe) & $|\Delta\text{DP}|$ & $|\Delta\text{EO}|$ & \textbf{DI} \\
\midrule
\textbf{Genre} & H 0.37, F 0.38   & 0.010 & 0.058 & 0.97 \\
\textbf{Âge}   & older 0.40, younger 0.27 & \textbf{0.125} & 0.066 & \textbf{0.69} \\
\bottomrule
\end{tabular}
\caption{Métriques d'équité du baseline.}
\end{table}

\begin{leconbox}[Constat]
Sur le \textbf{genre}, biais quasi nul après retrait de \texttt{personal\_status\_sex} (DI = 0.97).
Sur l'\textbf{âge}, $|\Delta\text{DP}| = 0.125$ persiste et DI = 0.69 viole la règle US 80\,\%. C'est le \textbf{point d'alerte principal} pour un audit de conformité.
\end{leconbox}

\subsection{Audit intersectionnel}
\label{sec:intersection}

L'audit séparé peut masquer des biais combinés. On crée un attribut à 4 niveaux \texttt{genre\_age} :

\begin{table}[H]
\centering
\begin{tabular}{lccc}
\toprule
\textbf{Sous-groupe} & \textbf{n} & \textbf{Approval rate} & \textbf{TPR} \\
\midrule
female\_older   & 36  & 0.44 & 0.85 \\
male\_older     & 120 & 0.38 & 0.85 \\
female\_younger & 25  & 0.28 & 0.92 \\
male\_younger   & 19  & \textbf{0.26} & 0.86 \\
\bottomrule
\end{tabular}
\caption{Taux d'octroi par sous-groupe intersectionnel.}
\end{table}

Synthèse comparative :

\begin{table}[H]
\centering
\begin{tabular}{lccc}
\toprule
\textbf{Axe} & $|\Delta\text{DP}|$ & $|\Delta\text{EO}|$ & \textbf{DI} \\
\midrule
Genre seul     & 0.010 & 0.058 & 0.97 \\
Âge seul       & 0.125 & 0.066 & 0.69 \\
\textbf{Intersectionnel} & \textbf{0.181} & \textbf{0.222} & \textbf{0.59} \\
\bottomrule
\end{tabular}
\caption{Écarts d'équité : agrégés vs intersectionnel.}
\end{table}

\begin{figure}[H]
\centering
\includegraphics[width=0.95\linewidth]{../outputs/14_intersectional.png}
\caption{Approval rate par sous-groupe (gauche, ligne rouge = seuil DI 80\,\% local) et écarts d'équité agrégés vs intersectionnel (droite).}
\end{figure}

\begin{leconbox}[Surprise]
Les \textbf{male\_younger} (n = 19) ont le taux d'octroi le plus bas (0.26), \textbf{pas} les female\_younger (n = 25, 0.28) comme l'intuition initiale le suggérait. L'intersection multiplie le biais par 1.5 par rapport à l'analyse par axe seul.
\end{leconbox}

% ════════════════════════════════════════════════════════════════
\section{Atténuation des biais}

\subsection{Pré-traitement : reweighing (Kamiran \& Calders 2012)}

On repondère chaque exemple pour rendre $S \perp Y$ dans la distribution pondérée :
$$w_i = \frac{P(S = s_i) \cdot P(Y = y_i)}{P(S = s_i, Y = y_i)}$$

Méthode model-agnostic : on ré-entraîne \texttt{LogisticRegression(sample\_weight=w)}, le code d'apprentissage est inchangé.

\subsection{Post-traitement : seuils par groupe (Hardt et al.\ 2016)}

Au lieu d'un seuil global $\tau$, on calibre $\tau_g$ par groupe sensible sur la validation pour égaliser le taux de sélection cible (\emph{Demographic Parity}). Aucun ré-entraînement requis.

\subsection{Comparaison des 4 configurations}
\label{sec:compare}

\begin{table}[H]
\centering
\small
\begin{tabular}{llcccc}
\toprule
\textbf{Attr.} & \textbf{Configuration} & \textbf{AUC} & \textbf{BalAcc} & $|\Delta\text{DP}|$ & $|\Delta\text{EO}|$ \\
\midrule
Genre & Baseline        & 0.785 & 0.669 & 0.010 & 0.058 \\
Genre & Reweighing      & 0.785 & 0.665 & 0.006 & 0.050 \\
Genre & Baseline + PP   & 0.785 & 0.669 & 0.105 & 0.196 \\
Genre & Reweighing + PP & 0.785 & 0.669 & 0.128 & 0.230 \\
\midrule
Âge   & Baseline        & 0.785 & 0.669 & 0.125 & 0.066 \\
Âge   & Reweighing      & 0.784 & 0.640 & 0.132 & 0.098 \\
Âge   & Baseline + PP   & 0.785 & 0.644 & \textbf{0.060} & \textbf{0.135} \\
Âge   & Reweighing + PP & 0.784 & 0.650 & 0.069 & 0.221 \\
\bottomrule
\end{tabular}
\caption{Performance et équité sur le test (PP = post-processing par seuils par groupe).}
\label{tab:compare}
\end{table}

\begin{figure}[H]
\centering
\includegraphics[width=0.95\linewidth]{../outputs/02_tradeoff.png}
\caption{Écarts d'équité (haut) et performance (bas) pour les 4 configurations sur les deux attributs sensibles.}
\end{figure}

\begin{attentionbox}[Théorème d'impossibilité \emph{en action}]
Sur l'âge, le post-processing réduit $|\Delta\text{DP}|$ de \textbf{0.125 à 0.060} mais \textbf{double} $|\Delta\text{EO}|$ (0.066 vers 0.135). On ne peut pas atteindre les deux objectifs simultanément.

Sur le genre, où le biais initial est négligeable, le post-processing introduit un biais artificiel : le seuil par groupe calibré sur 200 lignes de validation ne généralise pas au test.
\end{attentionbox}

% ════════════════════════════════════════════════════════════════
\section{Interprétabilité}

\subsection{SHAP global}

Pour la régression logistique, les valeurs de Shapley admettent une expression exacte (Lundberg \& Lee 2017, Corollary 1) :
$$\phi_i(x) = w_i \cdot (x_i - \mathbb{E}[x_i])$$

Implémentation via \texttt{shap.LinearExplainer}. L'importance est l'agrégation $\frac{1}{n}\sum_x |\phi_i(x)|$.

\begin{figure}[H]
\centering
\includegraphics[width=0.55\linewidth]{../outputs/03_shap.png}
\caption{Top 10 features par Mean $|$SHAP$|$ sur le test.}
\end{figure}

Le modèle s'appuie sur des variables \emph{financières légitimes} : \texttt{checking\_status} (0.67), \texttt{credit\_history} (0.49), \texttt{installment\_rate} (0.39), \texttt{savings\_account\_bonds} (0.36). Variables démographiques absentes du top.

\subsection{SHAP local : explication individuelle}

Décomposition de la prédiction d'un client unique. Utile pour répondre au \textbf{droit d'explication} (RGPD article 22).

\begin{figure}[H]
\centering
\includegraphics[width=0.95\linewidth]{../outputs/09_shap_local.png}
\caption{Contributions par feature pour deux clients : haut risque (gauche, score = 0.95) et bas risque (droite, score = 0.01).}
\end{figure}

Sur le client à haut risque (vrai défaut), les features qui poussent vers défaut sont \texttt{credit\_amount} (+1.18), \texttt{checking\_status} (+0.81), \texttt{duration\_in\_month} (+0.66). Le modèle peut donc justifier sa décision auprès du client.

\subsection{Détection de proxies}
\label{sec:proxies}

Question : est-ce que les autres features encodent indirectement les attributs sensibles ?

On entraîne une régression logistique \texttt{features $\to$ attribut sensible}. L'AUC mesure objectivement la prédictibilité de l'attribut depuis le reste.

\begin{table}[H]
\centering
\begin{tabular}{lcl}
\toprule
\textbf{Attribut sensible} & \textbf{AUC proxy} & \textbf{Interprétation} \\
\midrule
Genre & 0.69          & signal proxy modéré \\
\textbf{Âge} & \textbf{0.86} & \textbf{très fortement reconstructible} \\
\bottomrule
\end{tabular}
\caption{Régression multivariée features vers attribut sensible.}
\end{table}

\begin{leconbox}[Recommandation Charlotte Laclau]
Retirer \texttt{age\_in\_years} ne suffit pas à \og{} cacher \fg{} l'âge : il est encodé dans \texttt{housing}, \texttt{employment}, \texttt{job}\dots \textbf{Conserver l'attribut sensible} et l'utiliser explicitement comme pondération dans une mitigation in-processing serait théoriquement plus efficace (\emph{fairness through awareness}, Dwork et al.\ 2012).
\end{leconbox}

\subsection{Comparaison SHAP : avant vs après retrait}

On entraîne deux modèles, l'un avec les attributs sensibles, l'autre sans. On compare la redistribution de l'importance SHAP par feature.

\begin{figure}[H]
\centering
\includegraphics[width=0.95\linewidth]{../outputs/08_shap_compare.png}
\caption{Top 12 features SHAP avec (bleu) vs sans (orange) attributs sensibles, et $\Delta$ SHAP (droite). Les barres vertes du panneau droit identifient les proxies absorbant l'effet retiré.}
\end{figure}

Top features absorbant l'effet : \texttt{present\_employment\_since} (+0.036), \texttt{telephone} (+0.022), \texttt{credit\_history} (+0.020). Confirmation graphique de l'AUC proxy = 0.86.

\subsection{Performance avec vs sans attributs sensibles}

Mesure du coût en performance et du gain en équité du retrait :

\begin{table}[H]
\centering
\begin{tabular}{lcccc}
\toprule
\textbf{Modèle} & \textbf{AUC} & \textbf{BalAcc} & $|\Delta\text{DP}_{\text{age}}|$ & \textbf{DI}$_{\text{age}}$ \\
\midrule
Avec attr. sensibles            & 0.781 & 0.660 & 0.110 & 0.69 \\
Sans attr. sensibles (baseline) & 0.785 & 0.669 & 0.125 & 0.69 \\
\bottomrule
\end{tabular}
\caption{Comparaison performance et équité du retrait des attributs sensibles.}
\end{table}

\begin{figure}[H]
\centering
\includegraphics[width=0.85\linewidth]{../outputs/16_metrics_with_vs_without.png}
\caption{Comparaison performance (gauche) et équité âge (droite) avec vs sans attributs sensibles.}
\end{figure}

\begin{leconbox}[Le retrait est cosmétique, pas correctif]
Le coût en performance est négligeable ($\Delta$ AUC = 0.004) et l'équité ne s'améliore pas (DI âge reste à 0.69). Sans intervention active (reweighing, post-processing, fairness through awareness), retirer la colonne sensible ne corrige pas le biais.
\end{leconbox}

% ════════════════════════════════════════════════════════════════
\section{Quantification d'incertitude}

\subsection{Bootstrap}

\textbf{Choix méthodologique} : sur conseil de l'enseignante, la robustesse adversariale demandée par le sujet est remplacée par une quantification d'incertitude par bootstrap (cf.\ cours \textit{Robustness and Uncertainty}, G.~Franchi).

On entraîne $B = 200$ régressions logistiques sur des échantillons bootstrap (Efron 1979) du train. Pour chaque modèle, on calcule les métriques sur le test, puis on agrège par quantiles 2.5\,\%/97.5\,\%.

\begin{table}[H]
\centering
\begin{tabular}{lcccc}
\toprule
\textbf{Métrique} & \textbf{Moyenne} & $\sigma$ & \textbf{IC 2.5\,\%} & \textbf{IC 97.5\,\%} \\
\midrule
AUC                       & 0.764 & 0.019 & 0.723 & 0.795 \\
BalAcc                    & 0.682 & 0.023 & 0.637 & 0.726 \\
$|\Delta\text{DP}|$ genre & 0.035 & 0.028 & 0.001 & 0.105 \\
$|\Delta\text{EO}|$ genre & 0.089 & 0.044 & 0.017 & 0.176 \\
$|\Delta\text{DP}|$ âge   & 0.104 & 0.046 & 0.020 & 0.202 \\
$|\Delta\text{EO}|$ âge   & 0.086 & 0.048 & 0.018 & 0.194 \\
\bottomrule
\end{tabular}
\caption{Intervalles de confiance bootstrap ($B = 200$).}
\end{table}

\begin{figure}[H]
\centering
\includegraphics[width=0.95\linewidth]{../outputs/05_bootstrap.png}
\caption{Distributions bootstrap : AUC (gauche), équité genre (centre), équité âge (droite).}
\end{figure}

\begin{attentionbox}[Lecture critique]
L'IC 95\,\% de $|\Delta\text{DP}|$ âge va de \textbf{0.02 à 0.20}. Le biais empirique (0.125) est dans cet IC, mais l'incertitude est plus large que la valeur. Annoncer \og{} biais $|\Delta\text{DP}| = 0.125$ \fg{} sans préciser \og{} $\pm 0.10$ \fg{} serait trompeur.
\end{attentionbox}

\subsection{Intervalles de prédiction par individu}

Pour chaque ligne du test, on agrège les $B$ scores en moyenne et écart-type. Une forte variance signale un client \textbf{sur la frontière de décision}.

\begin{figure}[H]
\centering
\includegraphics[width=0.95\linewidth]{../outputs/06_indiv_uncertainty.png}
\caption{Intervalles de prédiction par individu (gauche) et incertitude vs score (droite). La forme parabolique confirme que l'incertitude est maximale au milieu (cas frontières).}
\end{figure}

Les 10 individus les plus incertains ont des écarts-type bootstrap de 0.20--0.24. Le candidat le plus incertain (id = 236) a un score moyen de 0.39 mais un IC de [0.03, 0.85] : selon le tirage du jeu d'entraînement, ce client serait soit accepté, soit refusé. Une expertise humaine est recommandée pour ces cas.

% ════════════════════════════════════════════════════════════════
\section{Tableau de bord récapitulatif}

\begin{figure}[H]
\centering
\includegraphics[width=0.95\linewidth]{../outputs/12_dashboard.png}
\caption{Synthèse des métriques clés du baseline. Ligne pointillée rouge = seuil DI 80\,\% (EEOC).}
\end{figure}

\begin{leconbox}[Synthèse]
\textbf{Point d'alerte unique} : DI âge = 0.69, sous le seuil légal US 0.8 (règle EEOC). C'est l'élément à signaler prioritairement dans un audit de conformité. Tout le reste (performance, équité genre, calibration) est dans les marges acceptables.
\end{leconbox}

% ════════════════════════════════════════════════════════════════
\section{Conclusion}
\label{sec:conclusion}

\subsection{Résultats principaux}

\begin{enumerate}[itemsep=0.4em]
    \item \textbf{Modèle.} Régression logistique L2, AUC = 0.785, BalAcc = 0.669. Convergence en 6 itérations.

    \item \textbf{Forte présence de proxies multivariés} (approche Charlotte). La régression \texttt{features $\to$ âge} atteint AUC = 0.86 ; \texttt{features $\to$ genre} atteint AUC = 0.69. Retirer les attributs sensibles ne les cache pas, leur effet se redistribue (vérifié visuellement par SHAP avant/après, voir \S5.4 et \S5.5 dans le notebook).

    \item \textbf{Discrimination sur l'âge plus marquée que sur le genre.} $|\Delta\text{DP}|$ âge = 0.125, DI = 0.69 (sous seuil légal US 0.8). Audit intersectionnel : $|\Delta\text{DP}| = 0.181$ (1.5$\times$ plus fort) ; les jeunes hommes (n = 19) sont les plus défavorisés.

    \item \textbf{Théorème d'impossibilité confirmé empiriquement.} PP réduit $|\Delta\text{DP}|$ âge (0.125 vers 0.060) mais double $|\Delta\text{EO}|$ (0.066 vers 0.135).

    \item \textbf{Significativité statistique faible.} Bootstrap : IC 95\,\% de $|\Delta\text{DP}|$ âge = $[0.02, 0.20]$. Sur 1000 lignes, les conclusions sur l'équité doivent être nuancées par cette incertitude.

    \item \textbf{Le retrait des attributs n'est pas une mitigation.} AUC quasi inchangée (0.785 sans vs 0.781 avec) et DI âge reste à 0.69 dans les deux cas. Une intervention active est nécessaire.
\end{enumerate}

\subsection{Choix méthodologiques}

\begin{itemize}
    \item Robustesse adversariale du sujet remplacée par \textbf{bootstrap} (conseil de l'enseignante).
    \item Coût métier UCI ($5 \cdot \text{FN} + 1 \cdot \text{FP}$) pour ancrer le choix de seuil dans la réalité bancaire.
\end{itemize}

\subsection{Limites}

\begin{itemize}
    \item \textbf{Petit dataset} (1000 lignes, test 200) : IC 95\,\% des écarts d'équité dépassent 0.15.
    \item \textbf{Approche linéaire} : capacité limitée à capturer des interactions non linéaires.
    \item \textbf{Reweighing simple} : on ne traite qu'une dépendance binaire $S \perp Y$. Des méthodes plus avancées (\texttt{fairlearn.ExponentiatedGradient}, transport optimal) pourraient mieux gérer les contraintes.
    \item \textbf{Pas de \emph{fairness through awareness} implémentée}, malgré l'AUC proxy = 0.86 qui la justifierait. Piste pour une suite.
\end{itemize}

\subsection{Outils d'assistance utilisés}

Conformément à l'exigence du sujet, des LLM (Claude, GPT) ont été utilisés pour : aide à la rédaction du code (pipeline sklearn, intégration fairlearn/shap), structuration du notebook et du rapport. Les décisions méthodologiques (choix des métriques, interprétation, conception du bootstrap) ont été prises et discutées en équipe.

% ════════════════════════════════════════════════════════════════
\section*{Références}

{\footnotesize
\begin{thebibliography}{99}
\itemsep=0.1em
\bibitem{chouldechova2017} A.~Chouldechova, \textit{Fair prediction with disparate impact}, Big Data 5(2), 2017. \href{https://arxiv.org/abs/1703.00056}{arXiv:1703.00056}.
\bibitem{kleinberg2017} J.~Kleinberg, S.~Mullainathan, M.~Raghavan, \textit{Inherent trade-offs in the fair determination of risk scores}, ITCS 2017. \href{https://arxiv.org/abs/1609.05807}{arXiv:1609.05807}.
\bibitem{hardt2016} M.~Hardt, E.~Price, N.~Srebro, \textit{Equality of Opportunity in Supervised Learning}, NeurIPS 2016. \href{https://arxiv.org/abs/1610.02413}{arXiv:1610.02413}.
\bibitem{kamiran2012} F.~Kamiran, T.~Calders, \textit{Data preprocessing techniques for classification without discrimination}, KAIS 33(1), 2012. \href{https://doi.org/10.1007/s10115-011-0463-8}{doi:10.1007/s10115-011-0463-8}.
\bibitem{dwork2012} C.~Dwork et al., \textit{Fairness through awareness}, ITCS 2012. \href{https://arxiv.org/abs/1104.3913}{arXiv:1104.3913}.
\bibitem{bellamy2018} R.~K.~E.~Bellamy et al., \textit{AI Fairness 360}, IBM, 2018. \href{https://arxiv.org/abs/1810.01943}{arXiv:1810.01943}.
\bibitem{bird2020} S.~Bird et al., \textit{Fairlearn: a toolkit for assessing and improving fairness in AI}, Microsoft 2020. \href{https://fairlearn.org/}{fairlearn.org}.
\bibitem{lundberg2017} S.~Lundberg, S.-I.~Lee, \textit{A Unified Approach to Interpreting Model Predictions}, NeurIPS 2017. \href{https://arxiv.org/abs/1705.07874}{arXiv:1705.07874}. Linear SHAP : Corollary 1.
\bibitem{efron1979} B.~Efron, \textit{Bootstrap methods}, Annals of Statistics 7(1), 1979. \href{https://doi.org/10.1214/aos/1176344552}{doi:10.1214/aos/1176344552}.
\bibitem{lakshminarayanan2017} B.~Lakshminarayanan et al., \textit{Deep Ensembles}, NeurIPS 2017. \href{https://arxiv.org/abs/1612.01474}{arXiv:1612.01474}.
\bibitem{pedregosa2011} F.~Pedregosa et al., \textit{Scikit-learn: Machine Learning in Python}, JMLR 12, 2011. \href{https://www.jmlr.org/papers/v12/pedregosa11a.html}{jmlr.org}.
\bibitem{cours} Cours IA708, F.~d'Alché-Buc, C.~Laclau, G.~Franchi, Télécom Paris.
\end{thebibliography}
}

\end{document}
